\section{Three Scaling Axes}
\label{sec:scaling}

\Cref{sec:architecture} defined the training-to-serving path: a rollout is generated by a resident base plus an adapter revision, training updates that adapter, and export returns a PEFT adapter file in serving layout. Scaling MinT stresses three parts of that path: large-base placement, exported adapter bytes, and adapter cache state. \textbf{Scale up} keeps LoRA RL usable when dense or MoE bases require model-parallel placement and sparse-routing metadata. \textbf{Scale down} makes the adapter revision, rather than a merged checkpoint, the object that crosses from training to serving. \textbf{Scale out} keeps a large addressable policy catalog selectable while each engine admits only bounded CPU-resident and GPU-active working sets.

\subsection{Scale Up: LoRA RL on Large Dense and MoE Bases}

\paragraph{Megatron placement.}
MinT uses Megatron training groups when the base model is too large for a single PEFT worker. Tensor parallelism shards dense tensors. Expert parallelism assigns MoE experts to rank groups. Dense-module LoRA tensors follow the tensor-parallel shards for the dense weights they modify. Per-expert LoRA tensors are keyed by expert id: the EP shard owns the experts assigned to it, and TP shards the expert matrices inside that owner group when the run also uses tensor parallelism. Shared-expert LoRA is stored once per EP shard and deduplicated during export.

\paragraph{Policy switching.}
The base shards stay resident across policies. A policy switch restores the adapter tensors, optimizer moments, and accumulated gradients on the ranks that will execute the next update. This keeps the policy training state small relative to the resident base while preserving Megatron's sharding rules.

\paragraph{Distributed export.}
Serving needs one PEFT adapter revision that can attach to a resident vLLM base deployment. Export converts the Megatron training view into that serving view: tensor-parallel LoRA slices are gathered, replicated tensors are deduplicated, and expert LoRA tensors are collected from their expert-parallel owners. The exported adapter is the policy revision that later appears in serving, evaluation, rollback, and catalog records.

\paragraph{MoE router replay.}
MoE RL requires training-time scoring to use the expert path that generated each rollout token. R3 identifies router mismatch as a concrete source of instability in MoE RL: small implementation or precision differences can route the same token to different experts during rollout and training~\citep{r3_moe_router2025,chiang2026routerreplay}. R3 records MoE expert routing. For Qwen3-style MoE runs, MinT stores selected expert ids with rollout records. Training can replay a recorded route when the backend can reconstruct that expert path. Training masks a token when the rollout record lacks selected expert ids or the training backend cannot map those ids to its expert-parallel layout.

\paragraph{Sparse-attention provenance.}
Dynamic sparse attention has a separate mismatch channel. In GLM-5 and GLM-5.1, the DSA indexer and top-$k$ path decide which tokens participate in sparse attention; small numerical differences can change that token set~\citep{glm5_2026,stevenchiang2026supportglm5inmint}. MinT removes observed implementation mismatches where the stack exposes a concrete cause: indexer RoPE layout, normalized query/key inputs, deterministic top-$k$ behavior, frozen indexer defaults, long-context THD/CP support, and LoRA loading for DSA target modules. Probability mismatch can remain after those fixes, so MinT uses IcePop-style rollout correction~\citep{ling_every_step2025}: when the training/rollout probability ratio leaves the configured lower--upper trusted band, the token receives zero importance weight. This mitigation filters unsafe scoring terms. It does not replay every DSA indexer choice, and it does not prove that training used the exact sparse-attention token set selected by the inference engine.

\begin{table}[H]
\centering
\scriptsize
\setlength{\tabcolsep}{4pt}
\renewcommand{\arraystretch}{1.20}
\caption{Representative model-family support validated by the current MinT stack.}
\label{tab:supported_model_families}
\fittowidth{%
\begin{tabular}{@{}M{0.16\textwidth}M{0.18\textwidth}M{0.28\textwidth}M{0.32\textwidth}@{}}
\toprule
\apphead Family & Model structure & Scale validated & MinT support path \\
\midrule
\appkey{Qwen3 series} & Dense and MoE variants & 0.6B/4B dense adapters; 30B-A3B and 235B-A22B MoE runs & Single-worker PEFT, Megatron MoE training, expert-route records, adapter export, and shared-base serving \citep{qwen3_2025,chiang2026routerreplay} \\
\addlinespace[3pt]
\appkey{Moonlight \& Kimi K2} & MLA MoE & Moonlight-16B-A3B bring-up; Kimi K2 1.04T countdown-task LoRA RL & MLA/MoE adapter placement, Megatron-Bridge conversion, vLLM serving, and trillion-parameter-class LoRA RL \citep{chiang2026routerreplay,liu2025Build,kimi_k2_2025} \\
\addlinespace[3pt]
\appkey{GLM-5 / GLM-5.1} & MLA, DSA, MTP, and MoE & Frontier agentic family with DSA/MTP training and serving bring-up & DSA/MTP training patches, DSA LoRA target mapping, vLLM custom-forward LoRA loading, bridge conversion, and IcePop rollout correction \citep{glm5_2026,stevenchiang2026supportglm5inmint,ling_every_step2025} \\
\bottomrule
\end{tabular}%
}
\end{table}


\subsection{Scale Down: Adapter-Only Training-Serving Handoff}

\paragraph{Adapter handoff bytes.}
The training-serving handoff uses the LoRA adapter as the checkpoint. A measured Qwen3-4B rank-32 PEFT adapter file is 264{,}310{,}274 bytes, about 252 MiB. Adapter byte size is determined by rank, target modules, dtype, and tensor layout. A four-billion-parameter bf16 base checkpoint has an 8.0 GB weight floor before metadata and optimizer state. The measured adapter is about 3.3\% of that base-weight floor, and the serving path can load it without materializing another full base copy. Since LoRA tensor bytes scale approximately linearly with rank for a fixed target-module set, the same target pattern at rank 1 would be about 7.9 MiB, or roughly 0.10\% of the bf16 base-weight floor. This supports the abstract's conservative 1\% claim: compact rank-1 settings can fall below 1\%, while broader target-module choices and metadata can increase the footprint. The system property remains the same: the crossing artifact is an adapter revision, not a full checkpoint.

\paragraph{Serving-compatible export.}
MoE export applies the same adapter-revision rule to sharded training runs. The trainer writes a PEFT adapter file in the tensor layout expected by vLLM, with optimizer state and rank-local training files removed. The file that crosses from training to serving is the adapter revision.

\paragraph{Served-score attribution.}
A served score names the exported adapter revision together with its compatible base model, target modules, prompt renderer, scorer, and serving path. When a row reports a served score, MinT compares isolated adapter loading against shared-base serving before attributing the served score to the exported revision. This check separates the policy result from loader and routing differences.

\subsection{Scale Out: Policy-Population Serving}
\label{sec:scale_out_residency}

\paragraph{From stored adapters to live policies.}
Scale-out begins when exported adapters stop being a small set of training files and become a large set of deployable policies. The catalog can grow with tenants, product variants, rollback points, personalization branches, evaluation snapshots, and research sweeps. A user-facing name resolves to an exported adapter revision, the revision resolves to a durable adapter file, and a serving actor maps that file to a local adapter id when the policy enters that actor. \Cref{tab:scale_out_cache_tiers} names the three tiers that a policy revision can occupy at once; the paragraphs below visit them in turn.

\begin{table}[t]
\centering
\scriptsize
\setlength{\tabcolsep}{3.2pt}
\renewcommand{\arraystretch}{1.10}
\caption{Three cache tiers separate addressability from local residency and same-batch execution. A policy revision can occupy any subset of these tiers at a given moment; the scales, lifetimes, and control surfaces differ.}
\label{tab:scale_out_cache_tiers}
\fittowidth{%
\begin{tabular}{@{}M{0.16\textwidth}M{0.15\textwidth}M{0.15\textwidth}M{0.22\textwidth}M{0.26\textwidth}@{}}
\toprule
\apphead Tier & Scale & Lifetime & Promotion / eviction & Measured in this paper \\
\midrule
\appkey{Addressable catalog} & $10^3$--$10^6$ entries & Durable (control plane) & Promoted by adapter export; retired manually & Catalog sweep 1k--100k (\cref{tab:e4_serving_summary}); $10^6$-entry sizing in Appendix~\cref{tab:app_fleet_model} \\
\addlinespace[3pt]
\appkey{CPU adapter cache} & Hundreds per engine & Per actor run & Promoted by router or cache-miss load; LRU under memory pressure & 369 / 550 cached on one engine (\cref{tab:e4_serving_summary}; Appendix~\cref{tab:app_cache_ladders}) \\
\addlinespace[3pt]
\appkey{GPU batch} & $\le 64$ distinct adapters & One decoding step & Promoted by the batch scheduler; released at end of step & Same-batch frontier $N{=}64$ (\cref{tab:e4_serving_summary}; Appendix~\cref{tab:app_cache_ladders}) \\
\bottomrule
\end{tabular}%
}
\end{table}

\paragraph{Catalog size is the request-name scale.}
The measured 1k--100k catalog sweep keeps name resolution separate from engine-local cache state. The Qwen3-30B TP=4 experiments resolve request names against catalogs up to 100k entries, keep hundreds of adapters cached near one engine, and run clean same-batch probes through 64 distinct adapters. Appendix~\cref{tab:app_fleet_model} extrapolates the measured single-engine limits to a $10^6$-entry accumulated catalog under a fleet-level routing model. These measurements split scale-out into three quantities: how many policy revisions can be named, how many can stay cached near an engine, and how many can execute in the current batch. The $10^6$ number is addressability, not simultaneous GPU residency.

\paragraph{Traffic locality drives cache state.}
Catalog membership and local cache state change on different time scales. Catalog entries live in the control plane. Cache entries are local to a serving actor. The same adapter revision can be addressable in the catalog, cached in CPU memory on one actor, active in the current GPU batch, evicted back to shared storage, and later loaded again. Recurring traffic should stay near an engine that already holds the selected policy, while weak-locality sweeps and rollout waves expose cache misses.

\paragraph{Cold loading is service work.}
A cache miss does more than read a small LoRA file. The serving actor fetches tensors, materializes loader objects, registers the adapter with the engine, and activates it before decoding can start. Different missing policies serialize through the engine load path, so latency grows as a staircase when many unique policies arrive together. Concurrent cache-miss requests for the same missing policy can share one load; requests for different missing policies remain separate load jobs. MinT therefore treats cold loading as scheduled service work with deduplication and bounded backpressure.

\paragraph{Serving contract.}
Policy-population serving requires controls after name lookup. Routing preserves CPU-cache reuse when traffic has locality. Batch construction respects the smaller same-batch adapter limit; the addressable catalog is larger than any realized batch. Cold loading is observable, retryable, and backpressured because it turns a stored adapter revision into an engine-local adapter object. The policy record and exported revision stay stable while the adapter moves among shared storage, the CPU cache, and the GPU batch.

\paragraph{Representation controls the load staircase.}
MoE LoRA makes the cold path visible because a moderate-size adapter can expand into tens of thousands of tiny tensor objects and hundreds of megabytes of cached state across TP workers. MinT packs these tensors into a serving representation with nearly unchanged declared bytes, reducing object fanout before engine registration. The detailed serving measurements in \cref{sec:e4_serving} show 37{,}248 tensors reduced to 672 and an 8.5--8.7$\times$ live-load speedup, with the packed live engine-load slice below 0.2 seconds in the measured bursts. This live-load number is one component of a longer cold path that also includes routing, queueing, fetch, and generation. The exported adapter remains the policy unit, while cold loading becomes an explicit adapter lifecycle stage that can be routed, prewarmed, throttled, and optimized.
